\documentclass[../main.tex]{subfiles}

\graphicspath{{../image/}{../figure/}}

\begin{document}

\chapter{容许波扇与 Riemann 问题}\label{chapter:ch09-容许波扇与riemann问题}

守恒律方程组在时空坐标均匀拉伸下保持不变这一性质,诱导出自相似解的存在;自相似解沿从时空某一焦点发出的直线射线保持常值。这类解刻画了收敛于焦点、并在那里相互作用产生一个跳跃间断的一族波,而该跳跃间断又被分解为一个向外的波扇。

本章研究著名的 Riemann 问题,其目标是将跳跃间断分解为波扇。解将以三种不同方式构造出来,即:(a) 经典方法,将前面各章遇到的初等中心解拼接在一起,也就是常值状态、连接常值状态的激波,以及由常值状态或接触间断为边界的中心稀疏波;(b) 使向外波扇的总熵产生极小化;(c) 消失粘性方法,它使用随时间变化的粘性,使得所得到的耗散方程组与原双曲方程组一样在坐标拉伸下不变。在此过程中将出现一种新型间断,称为\emph{delta 激波}。

波扇的容许性问题将被提出。特别地,将考察以上述任一方法构造的解中所含的激波是否必定容许。

接下来,将确定最佳地逼近由两个波扇相互作用所生成的复杂波型的那个波扇。

我们还将展示一个方程组,其有界初值会产生共振波型,在有限时间内把解的振幅驱至无穷。

\subfile{Chapter-09/section-01.tex}

\subfile{Chapter-09/section-02.tex}

\subfile{Chapter-09/section-03.tex}

\subfile{Chapter-09/section-04.tex}

\subfile{Chapter-09/section-05.tex}

\subfile{Chapter-09/section-06.tex}

\subfile{Chapter-09/section-07.tex}

\subfile{Chapter-09/section-08.tex}

\subfile{Chapter-09/section-09.tex}

\subfile{Chapter-09/section-10.tex}

\subfile{Chapter-09/section-11.tex}

\end{document}
